\documentclass[../../../../uem_mat_mei_q.tex]{subfiles}

\begin{document}
    連立不等式

    \myspaceb%
    $\displaystyle%
        \left\{
        \begin{myaligna}
            x+2y-3 & \leqq 0 \\
            3x-7y-9 & \leqq 0 \\
            -6x+y-21 & \leqq 0
        \end{myaligna}%
        \right.
    $\\
    により表される座標平面上の領域を$A$とする．点$(x,\;y)$が領域$A$上を動くとき，%
    以下の問いに答えよ．

    \begin{myenub}
        \item 領域$A$の面積は$\myfracc{\myboxe{アイ}}%
            {\myboxe{ウ}}$である．
        \item $x^2-2x+y^2-4y$ の最大値は\:\myboxe{エオ}\:である．
        \item $x^2-2x+y^2-4y$ は

            \myspaceb%
            $\displaystyle%
                (x,\;y)=\left(\myfracc{\myboxe{カ}}{\myboxe{キ}},\;%
                \myfracc{\myboxe{ク}}{\myboxe{ケ}}\right)%
            $\\
            のとき最小値 $-\myfracc{\myboxe{コサ}}{\myboxe{シ}}$ をとる．
    \end{myenub}
\end{document}
